\DocumentMetadata{
  pdfversion=2.0,
  pdfstandard=ua-2,
  lang=en-US,
  tagging=on,
  tagging-setup={math/setup=mathml-SE,tabsorder=structure}
}
\documentclass[11pt,letterpaper]{article}
\usepackage[margin=0.68in,includefoot]{geometry}
\usepackage{newcomputermodern}
\usepackage{amsmath,amscd,mathtools}
\usepackage{tikz,adjustbox}
\providecommand{\square}{\mdlgwhtsquare}
\usepackage[hidelinks]{hyperref}
\hypersetup{
  pdftitle={Algebra First-Year Exam, August 2015, Part 2},
  pdfauthor={University of Florida Department of Mathematics},
  pdfsubject={Graduate mathematics examination},
  pdfdisplaydoctitle=true,
  bookmarksopen=true
}
\setcounter{secnumdepth}{0}
\setlength{\parindent}{0pt}
\setlength{\parskip}{0.28em}
\setlength{\emergencystretch}{3em}
\setlength{\leftmargini}{2.15em}
\setlength{\leftmarginii}{2.65em}
\setlength{\leftmarginiii}{3em}
\pagestyle{empty}
\raggedbottom
\allowdisplaybreaks
\NewTaggingSocketPlug{block/list/label}{examproblem}{%
  \tagstructbegin{tag=Lbl}\tagstructend%
  \tagstructbegin{tag=\LBody}%
  \tagstructbegin{tag=\ProblemHeadingTag,title-o={\ProblemHeadingText}}%
  \tagmcbegin{tag=\ProblemHeadingTag,actualtext={\ProblemHeadingText}}%
  #2%
  \tagmcend%
  \tagstructend%
}
\newenvironment{examproblems}[2]{%
  \def\ProblemHeadingTag{#2}%
  \AssignTaggingSocketPlug{block/list/label}{examproblem}%
  \begin{enumerate}%
  \setcounter{enumi}{#1}%
  \renewcommand{\labelenumi}{\textbf{\theenumi.}}%
  \setlength{\topsep}{0.35em}%
  \setlength{\partopsep}{0pt}%
  \setlength{\itemsep}{0.48em}%
  \setlength{\parsep}{0.18em}%
}{\end{enumerate}}
\newenvironment{examsubparts}{%
  \AssignTaggingSocketPlug{block/list/label}{default}%
  \begin{enumerate}%
  \setlength{\topsep}{0.18em}%
  \setlength{\partopsep}{0pt}%
  \setlength{\itemsep}{0.16em}%
  \setlength{\parsep}{0.1em}%
}{\end{enumerate}}
\newenvironment{examinstructions}{%
  \par\begingroup\itshape
}{%
  \par\endgroup\vspace{0.18em}
}
\makeatletter
\renewcommand\subsection{\@startsection{subsection}{2}{0pt}%
  {-1.25ex plus -.25ex minus -.1ex}%
  {0.55ex plus .1ex}%
  {\normalfont\large\bfseries}}
\renewcommand{\@maketitle}{%
  \newpage\null\vskip -1em%
  {\footnotesize\bfseries
    \href{https://www.ufl.edu/}{University of Florida}, \href{https://math.ufl.edu/}{Department of Mathematics}\par}%
  \vskip -0.5em%
  \tagstructbegin{tag=H1,title={Algebra First-Year Exam, August 2015, Part 2}}%
  \tagpdfparaOff%
  \tagmcbegin{tag=H1}%
    {\LARGE\bfseries\href{https://gma.math.ufl.edu/exams/algebra-first-year/}{Algebra First-Year Exam}, August 2015, Part 2\par}%
  \tagmcend%
  \tagpdfparaOn%
  \tagstructend%
  \vskip 0.25em%
  \tagmcbegin{artifact}%
    \hrule height 1pt%
  \tagmcend%
  \vskip 1em%
}
\makeatother
\title{Algebra First-Year Exam, August 2015, Part 2}
\author{}
\date{}
\begin{document}
\maketitle
\thispagestyle{empty}
\begin{examinstructions}
Answer four problems. (If you turn in more, the first four will be graded.) Within reason, you may quote theorems as long as you state them clearly.
\end{examinstructions}
\begin{examproblems}{0}{H2}
\def\ProblemHeadingText{Problem 1.}
\item
Let~\(R\) be a principal ideal domain. Prove that it satisfies the ascending chain condition on ideals.
\def\ProblemHeadingText{Problem 2.}
\item
This problem has three parts.
\begin{examsubparts}
\item[\textnormal{(a)}]
Define Unique Factorization Domain.
\item[\textnormal{(b)}]
Define irreducible element of an integral domain~\(R\).
\item[\textnormal{(c)}]
Prove from your definitions, that if~\(R\) is a Unique Factorization Domain and~\(p\) is an irreducible element of~\(R\), then the ideal \(P=(p)\) generated by~\(p\) is a prime ideal of~\(R\).
\end{examsubparts}
\def\ProblemHeadingText{Problem 3.}
\item
This problem has two parts.
\begin{examsubparts}
\item[\textnormal{(a)}]
Define Euclidean domain.
\item[\textnormal{(b)}]
Prove that every Euclidean domain is a principal ideal domain.
\end{examsubparts}
\def\ProblemHeadingText{Problem 4.}
\item
Let~\(R\) be an integral domain and let~\(M\) be an~\(R\)-module. Suppose that~\(N\) is an~\(R\)-submodule of~\(M\). We denote by \(\operatorname{Ann}(M)\) the annihilator of~\(M\) in~\(R\), so that \(\operatorname{Ann}(M)\) is an ideal of~\(R\).
\begin{examsubparts}
\item[\textnormal{(a)}]
Prove that \(\operatorname{Ann}(M/N)\operatorname{Ann}(N)\subseteq\operatorname{Ann}(M)\).
\item[\textnormal{(b)}]
Give an example (with proof) where \(\operatorname{Ann}(M/N)\operatorname{Ann}(N)\ne\operatorname{Ann}(M)\).
\item[\textnormal{(c)}]
Give an example (with proof) where \(\operatorname{Ann}(M/N)\operatorname{Ann}(N)=\operatorname{Ann}(M)\).
\end{examsubparts}
\def\ProblemHeadingText{Problem 5.}
\item
How many distinct isomorphism types of abelian groups of order \(2^4\cdot 3^4\cdot 5^4\) and exponent \(2^3\cdot 3\cdot 5^2\) are there? Justify your answer carefully.
\def\ProblemHeadingText{Problem 6.}
\item
Let~\(K\) be a splitting field for \(x^4+2\) over~\(\mathbb{Q}\). Calculate the degree of \([K:\mathbb{Q}]\), justifying carefully each step.
\end{examproblems}
\end{document}
